\section{Operational Use of Temporal Validity}\label{sec:online_adaptation}

The temporal-validity profile takes operational form through three ingredients: a
finite-sample regime, a local drift geometry, and a criterion that relates
estimation uncertainty to the local validity curvature of the useful-memory region. The horizon
law supplies the target once the regime parameter $a$ and the local roughness
pair $(H,\zeta)$ are known. Lag-geometry information must be strong enough to
place the plug-in horizon inside the useful-memory region with high probability.

A direct plug-in rule based on a local slope and intercept is too unstable to be
useful. A structural controller first selects the finite-sample regime, then
estimates roughness from aggregated lag discrepancies, and finally maps those
quantities into an admissible useful-memory band. On the synthetic phase
profile, that controller improves substantially on the naive plug-in rule and on
a pure activity baseline. Adaptive EWMA and adaptive scalar Kalman baselines are
more competitive, especially on the smooth default profile, but they degrade
more visibly in rough and alternating regimes.

Across the tested default, smooth, rough, alternating, and ramp-up profiles with
three seeds and a bounded window-jump constraint, the structural controller
improves on the direct plug-in rule in every run, beats the activity baseline
and adaptive EWMA in $13/15$ runs each, and beats the adaptive scalar Kalman
baseline in $11/15$ runs. The smooth default profile remains the most
competitive setting for the classical filters, but the rough and alternating
profiles favor the band-constrained structural rule clearly. The tested profiles
place the structural controller ahead of the direct plug-in rule, the activity
baseline, adaptive EWMA, and the adaptive scalar Kalman baseline;
see \Cref{tab:online_controller_summary}.
